\documentclass{article}
\usepackage{graphicx} % Required for inserting images
\usepackage{amsmath}
\usepackage[english, russian] {babel}
\usepackage[utf8]{inputenc}
\usepackage[T2A]{fontenc}
\usepackage{minted}
\usepackage{float}
\usepackage{amssymb}

\title{Базы данных}
\author{silvia.lesnaia }


\begin{document}

\maketitle

\textbf{16.09.25}

\section{}

Стратиегии поддержиания ссылочнйчной целостности

Restrict(огрнаичить) - не разрешать выполнение операции, 
приводящей к нарушениию ссылочной целостности. Самая простая стратегия,
требующая только проверки, имеются ли только кортежи в дочернем отношении, связанные с некоторым 
кортежом в родительском отношении.


Cascade(каскадировать) - разрешить выполнение требуемой оперции, но внести при этом неопходимые правки
в других отношениях так, чтоб не допустить нарушения ссылочно целостности и сохрнаить имеющие связи.

Ссылочная целостность поддреживаться двумя способами:

 - На уровне определения данных 

 ТУТ КОД ДОЛЖЕН БЫТЬ, В ПРЕЗЕНТАЦИИ НЕПОНЯТНО ЧТО ЭТО ДАЖЕ НЕ SQL

 - С использованием триггеров

  ТУТ КОД ДОЛЖЕН БЫТЬ, В ПРЕЗЕНТАЦИИ НЕПОНЯТНО ЧТО ЭТО ДАЖЕ НЕ SQL

Триггер - специальная функциия, хранимая процедура которая запускается автоматически. 
Триггер можно октолючить, но для этого нужны права. А так работает он автоматом. Он автоматичски 
целпяется.

В триггер можно запихнуть номральные сообщния об ошибках для пользователей. Т.е он может 
слуэить неким обработчиком.

\section{Реляционная алгебра}

В описании реляционной модели утвежрдается, что доступ к реляцционым данным осущетсвлется при 
помощи реляционной алгберы(РА) или эквивалетоного ему реляционного исчиосления(РИ). В реализицих 
реляионных СУБД сейчас не тспользуется в чистом виде ни РА, ни РИ.

Стандарт доступа к реляционным данным стал язык SQL(Structed Query Language).


Язык доступа к данным должен быть реляционно полным.

Реляыионныз опертаоры:

Теоретико-множественные операторы:

тут лолжны быть списки 



а тут опаерции по типу объедеения или пересечения или вычитания или тд. 

\end{document}